\subsection*{الف}
سیگنال
$x(t)=\sin(2\pi\cdot 8t)$
دارای بیشینهٔ فرکانس
$f_{\max}=8~\text{Hz}$
است. شرط نایکوئیستِ امن برای نمونه‌برداری بدون آلیاسینگ:
\[
f_s \ge 2 f_{\max} = 16~\text{Hz}.
\]
در حالی که
$f_s=8.5~\text{Hz} < 16~\text{Hz}$،
پس شرط نایکوئیست برقرار نیست و آلیاسینگ رخ می‌دهد.

\subsection*{ب}
فرکانس گسسته‌زمانی:
\[
\omega_0 = 2\pi \frac{f_0}{f_s} = 2\pi\frac{8}{8.5}=2\pi\frac{16}{17}.
\]
با توجه به تناوب در حوزهٔ گسسته‌زمان، داریم:
\[
2\pi\frac{16}{17} \equiv -2\pi\frac{1}{17}\quad (\text{مد }2\pi),
\]
پس فرکانس ظاهری در \lr{baseband} برابر است با
\[
f_{\text{alias}}=\frac{1}{17}f_s = \frac{8.5}{17}=0.5~\text{Hz}.
\]

\subsection*{ج}
بین دو نمونهٔ متوالی، فاز سیگنالِ نمونه‌برداری‌شده به اندازهٔ
$\omega_0=2\pi\frac{16}{17}$
جلو می‌رود که معادلِ عقب‌رفتنِ کوچکِ
$2\pi\frac{1}{17}$
است. بنابراین دنبالهٔ نمونه‌ها مانند یک سینوس کم‌فرکانس دیده می‌شود، چون در زمان‌های نمونه‌برداری، سینوس پر‌فرکانس و یک سینوس کم‌فرکانس (پس از تا شدن طیف) مقادیر یکسانی تولید می‌کنند.

\subsection*{د}
یک سیگنال متفاوت که دقیقاً همان دنبالهٔ نمونه‌ها را بدهد:
\[
x_2(t)=\sin(2\pi\cdot 16.5\, t).
\]
زیرا برای $t=n/f_s$:
\begin{flalign*}
x_2[n]
&=\sin\!\left(2\pi\cdot 16.5 \frac{n}{8.5}\right)
=\sin\!\left(2\pi\cdot \frac{33}{17}n\right)
=\sin\!\left(2\pi\left(1+\frac{16}{17}\right)n\right) \\
&=\sin\!\left(2\pi n + 2\pi\frac{16}{17}n\right)
=\sin\!\left(2\pi\frac{16}{17}n\right)
=x[n]. &&
\end{flalign*}

\subsection*{هـ}
\[
f(x,y)=2\cos(2\pi(3x+4y))
= e^{j2\pi(3x+4y)} + e^{-j2\pi(3x+4y)}.
\]
پس طیف فوریهٔ دوبعدی شامل دو ضربه (دلتا) در نقاط پشتیبانی زیر است:
\[
F(u,v) = \delta(u-3,v-4) + \delta(u+3,v+4),
\]
یعنی نقاط
$(u,v)=(3,4)$
و
$(u,v)=(-3,-4)$.

\subsection*{و}
با
$\Delta x=\Delta y=0.2$
فرکانس‌های نمونه‌برداری:
\[
f_{sx}=\frac{1}{\Delta x}=5,\qquad
f_{sy}=\frac{1}{\Delta y}=5 \qquad (\text{چرخه بر واحد طول}).
\]
حد نایکوئیست در هر راستا:
\[
\frac{f_{sx}}{2}=2.5,\qquad \frac{f_{sy}}{2}=2.5.
\]
اما مؤلفه‌های طیفیِ تصویر در
$u=\pm 3$
و
$v=\pm 4$
قرار دارند که هر دو خارج از بازه‌های
$[-2.5,2.5]$
هستند؛ بنابراین شرط نایکوئیست در هر دو جهت نقض می‌شود و آلیاسینگ رخ می‌دهد.

\subsection*{ز}
نمونه‌برداری روی شبکهٔ مستطیلی، طیف را با دوره‌های
$(f_{sx},f_{sy})=(5,5)$
کپی می‌کند. مؤلفهٔ
$(3,4)$
با شیفتِ
$(-5,-5)$
به داخل ناحیهٔ عبور فیلتر بازسازی (مستطیل نایکوئیست) می‌افتد:
\[
(3,4)-(5,5)=(-2,-1).
\]
همچنین
$(-3,-4)$
با شیفتِ
$(+5,+5)$
به
\[
(-3,-4)+(5,5)=(2,1)
\]
منتقل می‌شود. پس خروجی بازسازی (با فیلتر پایین‌گذر ایده‌آل با برش‌های
$f_{sx}/2$
و
$f_{sy}/2$)
دارای دو ضربه در
$(2,1)$
و
$(-2,-1)$
خواهد بود، یعنی:
\[
\hat f(x,y)=2\cos\big(2\pi(2x+y)\big).
\]

\subsection*{ح}
اگر
$\Delta x=\Delta y=0.1$
باشد، آنگاه
\[
f_{sx}=f_{sy}=10,\qquad \frac{f_{sx}}{2}=\frac{f_{sy}}{2}=5.
\]
در این حالت نقاط طیفی
$(\pm 3,\pm 4)$
کاملاً داخل ناحیهٔ نایکوئیست
$[-5,5]\times[-5,5]$
قرار می‌گیرند؛ بنابراین آلیاسینگ رخ نمی‌دهد و بازسازی ایده‌آل دقیقاً همان تصویر اصلی را برمی‌گرداند:
\[
\hat f(x,y)=2\cos\big(2\pi(3x+4y)\big).
\]